In \cite{MR3694599}, Li characterized continuous orbit equivalence of topologically free partial group actions in terms of diagonal-preserving isomorphisms of the associated C*-crossed products. In this section, we will extend the notion of continuous orbit equivalence to partial actions of inverse semigroups and characterize orbit equivalence of topologically free systems in terms of diagonal-preserving isomorphisms of the associated skew inverse semigroup rings.

Recall that if $\theta=\left(\theta_s:X_{s^*}\to X_s\right)_{s\in S}$ is a partial action of an inverse semigroup $S$ on a topological space $X$, we write $S\ast X=\left\{(s,x)\in S\times X:x\in X_{s^*}\right\}$.

\begin{definition}\label{definitioncoe}\index{Continuous orbit equivalence}
Let $\theta=\left(\theta_s:X_{s^*}\to X_s\right)_{s\in S}$ and  $\gamma=\left(\gamma_t:Y_{t^*}\to Y_t\right)_{t\in T}$ be topological partial actions on topological spaces $X$ and $Y$, respectively. We say that $\theta$ and $\gamma$ are \emph{continuously orbit equivalent} if there exists a homeomorphism 
\[\varphi: \: X \longrightarrow Y\]
and continuous maps
\[a: S\ast X \longrightarrow T\qquad\text{and}\qquad b: T\ast Y \longrightarrow S\]
such that for all $x\in X$, $s\in S_x$, $y\in Y$ and $t\in T_y$,
\begin{enumerate}%[{\rm (i)}]\setlength\itemsep{1ex}
\item[(i)] $\varphi(\theta_s(x)) = \gamma_{a(s,x)}(\varphi(x))$;
\item[(ii)] $\varphi^{-1}(\gamma_t(y)) = \theta_{b(t,y)}(\varphi^{-1}(y))$.
\end{enumerate}
Implicitly, we require that $a(g,x) \in T_{\varphi(x)}$ and $b(t,y) \in S_{\varphi^{-1}(y)}$. We will call the triple $(\varphi,a,b)$ a \emph{continuous orbit equivalence} between $\theta$ and $\gamma$.
\end{definition}

%\begin{example}
%...
%\end{example}

Our next goal is to prove, under appropriate conditions, that continuous orbit equivalence of topologically free partial actions is equivalent to isomorphism of the respective groupoids. This is a generalization of the analogous results for groups (\cite[Theorem 1.2]{MR3789176} and \cite[Theorem 2.7]{MR3694599}).

Given a continuous orbit equivalence $(\varphi,a,b)$ between partial actions $\theta$ and $\gamma$ as in Definition \ref{definitioncoe}, we will factor the map $a:S\ltimes_\theta X\to T$ to a product-preserving map\ftnt{We refrain from using the word ``morphism'' in this case since groupoids and inverse semigroups do not belong to a common category that we have previously defined.} from the groupoid of germs $S\ltimes_\theta X$ to $T$, and similarly for the map $b:T\ltimes_\gamma Y\to S$. For this, we need to prove some identities related to how $(\varphi,a,b)$ respects the structure of $S$ and $T$.

\begin{lemma}\label{lem:seila}
Let $\theta=\left(\theta_s:X_{s^*}\to X_s\right)_{s\in S}$ and  $\gamma=\left(\gamma_t:Y_{t^*}\to Y_t\right)_{t\in T}$ be topologically free partial actions, and $(\varphi,a,b)$ is a continuous orbit equivalence between $\theta$ and $\gamma$. Assume that the groupoids of germs $S\ltimes_\theta X$ and $T\ltimes_\gamma Y$ are Hausdorff. Then
\begin{enumerate}%\setlength\itemsep{1ex}
\item[(a)] $[s_1, x]=[s_2, x]$ implies that $[a(s_1,x),\varphi(x)]=[a(s_2,x),\varphi(x)]$, for all $x\in X$ and $s_1, s_2 \in S_x$.

\item[(b)] $[a(s_1s_2, x),\varphi(x)]=[a(s_1, \theta_{s_2}(x))a(s_2, x),\varphi(x)]$ for all $x\in X$ and $s_2\in S_x$ and $s_1\in S_{\theta_{s_2}(x)}$.

\item[(c)] $[b(a(s, x), \varphi(x)),x] = [s,x],$ for all $x\in X$ and $s\in S_x$.
\end{enumerate}
Analogous statements hold with $(\varphi^{-1},b,a)$ in place of $(\varphi,a,b)$.
\end{lemma}

\begin{proof}
\begin{enumerate}
\item[(a)] Let $x\in X$ and $s_1, s_2 \in S_x$. Choose an open neighbourhood $U$ of $x\in X$ such that
\[a(s_1, \widetilde{x})=a(s_1, x)\qquad\text{and}\qquad a(s_2, \widetilde{x})=a(s_2, x)\]
whenever $\widetilde{x}\in U$. Then for all $\widetilde{x}\in U\cap\varphi^{-1}(\Lambda(\gamma))$,
\begin{align*}
[s_1,\widetilde{x}]=[s_2,\widetilde{x}]
      & \ \Longrightarrow \ \theta_{s_1}(\widetilde{x})=\theta_{s_2}(\widetilde{x})
        \ \Longrightarrow \ \varphi(\theta_{s_1}(\widetilde{x}))=\varphi(\theta_{s_2}(\widetilde{x}))  \\
      & \ \Longrightarrow \ \gamma_{a(s_1,\widetilde{x})}(\varphi(\widetilde{x}))=\gamma_{a(s_2,\widetilde{x})}(\varphi(\widetilde{x})).        
\end{align*}
so the definition of $\Lambda(\gamma)$ and the given property of $U$ imply $[a(s_1,x),\varphi(\widetilde{x})]=[a(s_2,x),\varphi(\widetilde{x})]$. Since $\gamma$ is topologically free, $\Lambda(\gamma)$ is dense in $Y$, so $U\cap\varphi^{-1}(\Lambda(\gamma))$ is dense in $U$. Taking the limit $\widetilde{x}\to x$ we conclude that $[a(s_1,x),\varphi(x)]=[a(s_2,x),\varphi(x)]$ (we are using the fact that $T\ltimes_\gamma Y$ is Hausdorff, so limits are unique).

\item[(b)] Choose an open neighbourhood $U$ of $x \in X$ such that 
\[a(s_1s_2,\widetilde{x})=a(s_1s_2,x),\quad a(s_1,\theta_{s_2}(\widetilde{x}))= a(s_1,\theta_{s_2}(x))\quad\text{and}\quad a(s_2, \widetilde{x}) = a(s_2,x)\]
for all $\widetilde{x} \in U$. Then for all $\widetilde{x} \in U \cap\varphi^{-1}(\Lambda(\gamma))$
\begin{align*}
  \gamma_{a(s_1s_2, \widetilde{x})}(\varphi(\widetilde{x})) 
     & = \varphi(\theta_{s_1s_2}(\widetilde{x})) 
       = \varphi(\theta_{s_1}(\theta_{s_2}(\widetilde{x})) 
       = \gamma_{a(s_1,\theta_{s_2}(\widetilde{x}))}(\varphi(\theta_{s_2}(\widetilde{x}))) \\
     & = \gamma_{a(s_1, \theta_{s_2}(\widetilde{x}))}(\gamma_{a(s_2, \widetilde{x})}(\varphi(\widetilde{x})))
       = \gamma_{a(s_1, \theta_{s_2}(\widetilde{x}))a(s_2, \widetilde{x})}(\varphi(\widetilde{x}))
\end{align*}
so, as in item (a), the given property of $U$ and the definition of $\Lambda(\gamma)$ imply
\[[a(s_1s_2,x),\varphi(\widetilde{x})]=[a(s_1, \theta_{s_2}(x))a(s_2, x),\varphi(\widetilde{x})].\]
Since $\varphi^{-1}(\Lambda(\gamma))\cap U$ is dense in $U$ we conclude that
\[[a(s_1s_2,x),\varphi(x)]=[a(s_1,\theta_{s_2}(x))a(s_2,x),\varphi(x)].\]

\item[(c)] Similarly to the previous items, take neighbourhoods $U$ of $x$ and $V$ of $\varphi(x)$ such that
\[a(s,\widetilde{x})=a(s,x)\qquad\text{and}\qquad b(a(s,x),\widetilde{y})=b(a(s,x),\varphi(x))\]
whenever $\widetilde{x}\in U$ and $\widetilde{y}\in V$. Then for all $\widetilde{x}\in U\cap\varphi^{-1}(V)\cap \Lambda(\theta)$,
\[\theta_{b(a(s,\widetilde{x}),\varphi(\widetilde{x}))}(\widetilde{x})=\varphi^{-1}(\gamma_{a(s,\widetilde{x})}(\varphi(\widetilde{x})))=\varphi^{-1}(\varphi(\theta_s(\widetilde{x})))=\theta_s(\widetilde{x})\]
so the properties of $U$, $V$ and $\Lambda(\theta)$ yield $[b(a(s,x),\varphi(x)),\widetilde{x}]=[s,\widetilde{x}]$ and again taking $\widetilde{x}\to x$ gives us the desired result.\qedhere
\end{enumerate}
%
%Items (d), (e) and (f) are analogous to (a), (b) and (c), by reversing the roles of $a$ and $b$, and of $\varphi$ and $\varphi^{-1}$.\qedhere
\end{proof}

\begin{theorem}\label{theo:coegroupoidgermiso}
Let $\theta=\left(\theta_s:X_{s^*}\to X_s\right)_{s\in S}$ and $\gamma=\left(\gamma_t:Y_{t^*}\to Y_t\right)_{t\in T}$ be topologically free, continuously orbit equivalent partial actions, and suppose that the groupoids of germs $S\ltimes_\theta X$ and $T\ltimes_\gamma Y$ are Hausdorff. Then $S \ltimes_\theta X$ and $T \ltimes_\gamma Y$ are isomorphic as topological groupoids.
\end{theorem}

\begin{proof}
Let $(\varphi,a,b)$ be a continuous orbit equivalence between $\theta$ and $\gamma$. By by Lemma~\ref{lem:seila}(a) we can define
\begin{align*}
\Phi:S \ltimes_\theta X  \longrightarrow T \ltimes_\gamma Y,\qquad \Phi([s,x])=[a(s,x), \varphi(x)]
\end{align*}
and by by Lemma~\ref{lem:seila}(b) it is a groupoid morphism. Since $a$ and $\varphi$ are continuous then $\Phi$ is continuous. Similarly, the map  
 \begin{align*}
\Psi:T \ltimes_\gamma Y & \longrightarrow S \ltimes_\theta X,\qquad\Psi([t,y])=[b(t,y), \varphi^{-1}(y)]
\end{align*}
is a continuous groupoid morphism as well. $\Phi$ and $\Psi$ are inverses of each other due to Lemma~\ref{lem:seila}(c).\qedhere
\end{proof}

Note that the continuous maps $a$ and $b$ in the definition of continuous orbit equivalence take values in discrete spaces (namely, the corresponding semigroups), and so $X$ and $Y$ are required to have sufficiently many clopen sets in order for a continuous orbit equivalence between the corresponding partial actions to exist. Since we will now be interested in constructing an orbit equivalence for two actions from an isomorphism of the corresponding groupoids of germs, we will need to concentrate on spaces which have sufficiently many clopen sets and partial actions which respect this structure.

\begin{definition}[{\cite[Definition~5.2]{MR2565546}}]
A topological partial action $\theta$ of an inverse semigroup $S$ on a topological space $X$ is \emph{ample} if
\begin{enumerate}
\item[(i)] $X$ is locally compact, Hausdorff and zero-dimensional;
\item[(ii)] $X_s$ is a compact-open subset of $X$ for all $s\in S$.
\end{enumerate}
\end{definition}

\begin{lemma}\label{lemmaisomorphismtocoe}
Let $\theta=\left(\theta_s:X_{s^*}\to X_s\right)_{s\in S}$ and $\gamma=\left(\gamma_t:Y_{t^*}\to Y_t\right)_{t\in T}$ be ample partial actions of inverse semigroups $S$ and $T$ on $X$ and $Y$, respectively, such that the groupoids of germs $S\ltimes_\theta X$ and $T\ltimes_\gamma Y$ are Hausdorff. Let $\Phi: S \ltimes_\theta X \rightarrow T \ltimes_\gamma Y$ be a topological isomorphism, and consider the restriction $\varphi:=\Phi|_X:X\to Y$. Then for all $s$ in $S$ there are elements $t_1, \ldots, t_n$ in $T$ and there are disjoint compact-open subsets $K_1, \ldots K_n$ of $Y$ such that:
\begin{enumerate}
\item[(a)] $K_i \subseteq Y_{{t_i}^*}$;
\item[(b)] $\varphi(X_{s^*})=\bigcup_{i=1}^nK_i$;
\item[(c)] $\left\{\varphi^{-1}(K_i):i=1,\ldots,n\right\}$ is a partition of $X_{s^*}$;
\item[(d)] For all $i\in\left\{1,\ldots,n\right\}$ and for all $x \in \varphi^{-1}(K_i)$, one has that $\Phi([s,x])=[t_i,\varphi(x)]$.
\end{enumerate}
\end{lemma}

\begin{proof}
Since $[s, X_{s^*}]$ is a compact-open bisection, then $\Phi([s, X_{s^*}])$ is a compact-open bisection in $T \ltimes_\gamma Y$, so there are elements $t_1, \ldots, t_n $ of $T$ and disjoint compact-open subsets $K_1,\ldots, K_n $ of $Y$ with $K_i \in Y_{t_i^*}$ such that
\[\Phi([s,X_{s^*}])=\bigcup_{i=1}^n[t_i,K_i].\ntag\label{equationlemmaisomorphismtocoeproof}\]

Then item (a) is trivially satisfied. Item (b) follows by taking sources on both sides of Equation (\ref{equationlemmaisomorphismtocoeproof}), and (c) follows from (b) and the fact that $\varphi$ is injective.

Let us prove (d). Consider $i\in\left\{1,\ldots,n\right\}$ and $x \in \varphi^{-1}(K_i)$. Then Equation (\ref{equationlemmaisomorphismtocoeproof}) implies that $\Phi[s,x]\in [t_j, K_j]$, for some $j \in \{i, \cdots, n \}$, and in particular $\so([s,x])\in K_j$. As $K_1,\ldots,K_n$ are pairwise disjoint and
\[\so(\Phi[s,x])=\Phi(\so[s,x])=\varphi(x)\in K_i,\]
then $K_j=K_i$. Therefore, $\Phi[s,x]=[t_i,\varphi(x)]$.\qedhere
\end{proof}

We are now ready to prove that topological isomorphisms between Hausdorff groupoids of germs yield a continuous orbit equivalence between the respective partial actions.

\begin{theorem}\label{theo:groupoidgermscoe}
Let $\theta=\left(\theta_s:X_{s^*}\to X_s\right)_{s\in S}$ and $\gamma=\left(\gamma_t:Y_{t^*}\to Y_t\right)_{t\in T}$ be ample partial actions of $S$ and $T$ on $X$ and $Y$, respectively, and suppose that the groupoids of germs $S \ltimes_\theta X$ and $T \ltimes_\gamma Y$ are isomorphic and Hausdorff. Then $\theta$ and $\gamma$ are continuously orbit equivalent.
\end{theorem}

\begin{proof}
Let $\Phi : S \ltimes_\theta X \rightarrow  T \ltimes_\gamma Y$ be an isomorphism of topological groupoids. Then 
\[\varphi:= \Phi|_X: X \rightarrow Y\]
is a homeomorphism.

Given $s\in S$, and choose $t_1,\ldots,t_n\in T$ and compact-open subsets $K_1$, \ldots, $K_n\subseteq Y$ satisfying properties (a)-(d) of Lemma~\ref{lemmaisomorphismtocoe}. Define $a(s,x)=t_i$ whenever $x\in\varphi^{-1}(K_i)$, so that $a$ is a continuous map on $\left\{s\right\}\times X_{s^*}$. This way, we define a continuous function $a$ on all of $S\ast X=\bigcup_{s\in S}\left\{s\right\}\times X_{s^*}$. Let us show that $a$ satisfies the desired property for a continuous orbit equivalence between $\theta$ and $\gamma$: Given $(s,x)\in S\ast X$, let $t=a(s,x)$, so the definition of $a(s,x)$ implies
\[\gamma_{a(s,x)}(\varphi(x))=\ra\left([a(s,x),\varphi(x)]\right)=\ra\left([t,\varphi(x)]\right)=\ra(\Phi[s,x])=\Phi(\ra[s,x])=\varphi(\theta_s(x)),\]
as desired.

Proceeding similarly with $\Phi^{-1}$ in place of $\Phi$, we construct a function $b:T\ast Y\to S$ with analogous properties, so that $a$ and $b$ describe a continuous orbit equivalence between $\theta$ and $\gamma$.\qedhere
\end{proof}

\begin{corollary}
Let $S$ be an inverse semigroup which is a weak semilattice and $\theta$ be an ample partial action of $S$ on $X$. Let $\tau$ be the canonical action of $\KB(S\ltimes_\theta X)$ on $X$ (see Example~\ref{examplecanonicalaction}). Then $\theta$ and $\gamma$ are continuously orbit equivalent.
\end{corollary}

Recall that in the definition of ample partial action we assume that the domains are compact-open. The conclusion of Theorem \ref{theo:groupoidgermscoe} is not true in general without this assumption, i.e., there are non-continuously orbit equivalent topologically free partial actions with isomorphic groupoids of germs.

\begin{example}\label{exampleomega1isbad}
Let $\omega_1$ be the first uncountable ordinal endowed with the order topology, making it zero-dimensional, locally compact and Hausdorff (see 1.19 and 17.2(c) of \cite{MR0264581}). Let $S=\left\{1\right\}$ be the trivial group, acting trivially on $\omega_1$.

Consider $T=\omega_1$ as a lattice, and the action $\gamma$ of $T$ on $\omega_1$ in which $\gamma_t$ is the identity function of $\left\{\alpha\in\omega_1:\alpha<t\right\}$ for all $t\in T$.

The transformation groupoids $S\ltimes\omega_1$ and $T\ltimes\omega_1$ are both unit groupoids and isomorphic to $\omega_1$. However, there is no continuous orbit equivalence between the trivial action of $S$ and $\gamma$.

Indeed, first we identify $S\ast\omega_1=\left\{1\right\}\times\omega_1$ with $\omega_1$ in the obvious manner. Let $\varphi:\omega_1\to\omega_1$ be a homeomorphism and $a:\omega_1\to T$ a continuous function. The condition in Definition \ref{definitioncoe} is that $\varphi(\alpha)$ belongs to the domain of $\gamma_{a(\alpha)}$ for all $\alpha$, that is,
\[\varphi(\alpha)<a(\alpha)\qquad\text{for all }\alpha\in\omega_1.\ntag\label{equationomega1isbad}\]
In other words, the continuous function $f:\omega_1\to\omega_1$, $f(\alpha)=a(\varphi^{-1}(\alpha))$ satisfies $\alpha<f(\alpha)$ for all $\alpha\in\omega_1$, so let us prove that this leads to a contradiction (the argument is similar to that of \cite[Example 20.11]{MR0264581}): Let $\alpha_0\in\omega_1$ be arbitrary and for all $n\geq 0$ set $\alpha_{n+1}=f(\alpha_n)$. Then $\beta=\sup_n\alpha_n$ satisfies $\beta=\lim_n\alpha_n$, and since $f$ is continuous then
\[f(\beta)=\lim_n f(\alpha_n)=\lim_n\alpha_{n+1}=\beta,\]
a contradiction.%Substituting $a(\alpha)$ by its successor in $\omega_1$, if necessary, we may assume that $\left\{a(\alpha)\right\}$ is clopen in $\omega_1$ for all $\alpha\in\omega_1$.
%
%
%Recall that $\omega_1$ is not paracompact. The usual proof of this fact is given by proving that $\left\{(-\infty,\beta):\beta\in\omega_1\right\}$ has no locally finite refinement. See \cite[Example 20.11]{MR0264581} for more details.
%
%However, if such functions $\varphi$ and $a$ exist consider, for each $\beta\in\omega_1$, the set
%\[a^{-1}(\beta)=\left\{\alpha\in\omega_1:a(\alpha)=\beta\right\}\]
%so $\left\{a^{-1}(\beta):\beta\in\omega_1\right\}$ is a clopen partition of $\omega_1$, and therefore
%\[\left\{\varphi(a^{-1}(\beta)):\beta\in\omega_1\right\}\]
%is also a clopen partition of $\omega_1$. However, if $\alpha\in\varphi(a^{-1}(\beta))$, then
%\[\varphi^{-1}(\alpha)\in a^{-1}(\beta),\quad\text{thus}\quad a(\varphi^{-1}(\alpha))=\beta\]
%and Equation (\ref{equationomega1isbad}) implies that $\alpha<\beta$. This means that $\left\{\varphi(a^{-1}(\beta)):\beta\in\omega_1\right\}$ is a clopen partition, and in particular locally finite, which refines
%\[\left\{(-\infty,\beta):\beta\in\omega_1\right\}\]
%which is impossible by Example 20.11 of \cite{MR0264581}.
\end{example}

Similar ideas are used in \cite[Example 20.11]{MR0264581} to prove that $\omega_1$ is not paracompact, which suggests that some property regarding paracompactness of the domains of the partial action needs to be satisfied in order for Theorem \ref{theo:groupoidgermscoe} to be valid. The general condition that allows the theorem to be proven is ultraparacompactness: a zero-dimensional, locally compact, Hausdorff space is \emph{ultraparacompact} if every open cover can be refined to a clopen partition. See \cite{MR0261565} and the references therein for distinct characterizations and more information about this property.

For example, every second-countable (or more generally every Lindelöf), zero-dimensional, locally compact Hausdorff space is ultraparacompact. The proof of (the analogue of) Theorem \ref{theo:groupoidgermscoe} with ultraparacompact domains instead of compact-open is essentially the same, but it is preferable to leave the theorem stated with more standard - and elegant - assumptions.

